\section{Structural Interpretation}


The previous section organized the effective region in which consciousness can obtain through speed difference, coupling viscosity, internal delay, effective torque-like quantity, and meaning gradient.

This section interprets structurally what those expressions mean.

The crucial point is that the expressions in the previous section do not generate consciousness.

They are not expressions for computing consciousness as an object; they are constraint conditions showing under what circumstances consciousness can open.

\subsection{Consciousness Is Not an Object}


In the position of this paper, consciousness is not an object existing internally.

It is not a component placed somewhere in the brain, content stored in memory, or the final output of computational processing.

Treating consciousness as an object immediately prompts questions: Where is it? Which circuit produces it? Which information measure does it correspond to? How much is generated?

But as shown in the previous section, consciousness is not reducible to a single variable.

Consciousness is not any one of $I$, $J$, $\Phi$, $\Delta v$, $\mu$, $\tau$, $\mathcal{T}_{\Phi}$.

It is the state that opens when the relation among these enters a certain effective region.

Consciousness is therefore not an object but a regime.

It is not a location but a state domain; not a product but a condition.

\subsection{Consciousness as a Window}


The consciousness window introduced in the previous section is a concept for grasping consciousness not as a point but as a window.

Within this window, the fast and slow layers do not fully synchronize.

Nor do they fully separate.

Processing does not flow away, yet does not become over-fixed.

Logs can sediment, but complete closure is not reached.

In this intermediate region, a referencing relation opens between current processing and past logs.

This is the conscious state as understood in this paper.

Consciousness is therefore not binary — ON or OFF.

Consciousness is a dynamic state maintained within the consciousness window, varying in stability, clarity, continuity, and directionality.

\subsection{Why Generation Is a Misleading Term}


When consciousness is said to be "generated," an implicit picture is at work: a generator exists, a product exists, and the generated thing exists somewhere.

This picture is valid for objects, representations, and outputs.

But in the framework of this paper, consciousness is not such an object.

Consciousness opens when the fast layer, slow layer, plasticity gate, meaning potential, speed difference, and coupling viscosity satisfy a certain relation.

This is less a matter of something new being created than of a referencing possibility opening between already-existing processing and logs.

Consciousness is therefore not modeled here as a generated object.

Consciousness opens.

This expression "opens" is not a poetic metaphor.

It is a structural description: the temporal passage for log referencing is temporarily established.

When the gate closes, consciousness does not persist — not because a generated object has been destroyed, but because the relation that made log referencing possible has been lost.

\subsection{Consciousness as Temporal Non-Closure}


From the VED perspective, consciousness is understood as temporal non-closure.

Complete synchrony eliminates speed difference.

Complete separation prevents flow from arriving.

Complete fixation eliminates plasticity.

Complete fluidization prevents log sedimentation.

Consciousness belongs to none of these extremes.

Consciousness is in a state that tends toward closure but does not close completely, tends to flow but does not flow away, tends to sediment but does not become over-fixed.

This state is called temporal non-closure.

Non-closure here is not incompleteness.

It is, rather, the condition for persistence.

If closure were complete, consciousness might be fixed as an object — but at that point it would no longer be a live referencing state.

Conversely, with no closure at all, processing would leave no trace and flow away.

Consciousness is between these two.

That is: consciousness is not the failure of complete closure but the stabilization of non-closure.

\subsection{Meaning Gradient and Selective Referencing}


From the IFGT perspective, consciousness is not merely a temporal structure.

Log referencing is directed by the meaning gradient.

In regions where information potential $\Phi$ is formed, processing is drawn toward specific logs along that gradient:

\begin{equation}
|\nabla \Phi| \uparrow \quad \Rightarrow \quad \text{referencing bias} \uparrow
\end{equation}

What appears in consciousness is therefore not random.

Strong memories, emotionally weighted events, repeatedly referenced concepts, and narratively organized logs tend to enter conscious referencing more easily.

However, a strong meaning gradient is not synonymous with strong consciousness.

An excessive meaning gradient fixes referencing to specific wells.

A meaning gradient that is too weak causes referencing to lose direction.

What consciousness requires is not the meaning gradient itself but a meaning gradient operating within the range of speed difference and viscosity.

In this respect, $\mathcal{T}_{\Phi}$ is important:

\begin{equation}
\mathcal{T}_{\Phi} \sim \kappa \frac{\Delta v}{\mu} |\nabla \Phi|
\end{equation}

This quantity does not determine alone what appears in consciousness.

But it is a structural index showing which logs are easily referenced and into which meaning wells processing tends to flow.

\subsection{Ego as a Derived Position}


In this structure, where does the ego stand?

In the position of this paper, the ego is not the subject of consciousness.

It is neither the owner of consciousness nor the center from which judgments are issued.

The ego is the position at which the speed difference between the fast and slow layers is transformed through Torque Converter Attention.

At this position, current processing and past logs intersect.

Action preparation, memory sedimentation, meaning gradient, and reportability simultaneously converge.

Humans tend to experience this intersection as the place where they are thinking, the place where they are choosing, the place where they are conscious.

But this does not mean the ego is the cause.

The ego is not a generator but a transformation point.

The ego is therefore not an illusion — it functions in a real sense.

But it is not a subject in a real sense either.

It is real as a transformation point of time-scale difference.

\subsection{Qualia as Boundary Signal}


Qualia, too, are repositioned within this structure.

Qualia are not the substance of consciousness.

They are the live foreground that appears at the boundary where fast-layer processing tends to sediment into the slow layer.

While remaining entirely in the fast layer, qualia have not yet become reportable.

Once fully sedimented into the slow layer, they become memory, label, report, concept.

Qualia are in between.

Qualia are therefore not stored.

What is stored is the log structure that remains after qualia have passed through.

In this view, qualia are not a mysterious entity.

But neither are they mere illusion.

They are the phenomenal surface of the boundary where log referencing is in the process of opening.

\subsection{Summary}


The structural interpretation of this section is as follows:

\begin{itemize}
\item Consciousness is not an object but a state domain.
\item Consciousness is not modeled as a generated object; it opens as log referencing.
\item Consciousness is maintained not as complete closure but as temporal non-closure.
\item What appears in consciousness is directed by the meaning gradient $\nabla\Phi$.
\item The ego is not a subject but the transformation point of speed difference.
\item Qualia are not the substance but the live foreground at the sedimentation boundary.
\end{itemize}

The structural interpretation of consciousness in this paper is therefore condensed in one sentence:

Consciousness is the state in which log referencing directed by the meaning gradient is open within a temporal non-closure regime that reaches neither complete synchrony nor complete separation.

The next section organizes how this structure actually operates, as the dynamics of log referencing.
